According to Kepler's second law we have;

The area is swept out at constant rate throughout the orbital motion.
Now, for this comet the area swept out in one orbital period is $\pi ab$ where
$a$ is the semi-major axis and $b$ the semi-minor axis of the orbit.
\begin{align*}
  a &= 16.0\ \text{A.U.} \\
  b^2 &= a^2 - (\text{distance from a focus to centre of ellipse})^2 \\
  &= (16.0)^2 - (16.0 - 0.5)^2 \\
  b &= 3.968\ \text{A.U.} \\
  \pi ab &= 199.5\ (\text{A.U.})^2
\end{align*}
Hence, the area swept out per year is
\[
  \frac{199.5\ (\text{A.U.})^2}{64\ \text{years}}
  = 3.1\ (\text{A.U.})^2/\text{year}
\]
